% GENERATED FILE. Edit canonical lessons, then run npm run generate:books.
% canonical-source-hash: fnv1a64:dbe90ab7718de814
% canonical-lessons: ML-C451-itapetuka, ML-C451-nirbandhikkuka, ML-C451-rekhappetuttuka, ML-C451-bhisanippetuttuka, ML-C451-matikkuka

\chapter{Pressure and Pushback}
\label{ch:ml-pressure-and-pushback}
\hlchaptermodality{451}

\begin{chapteropening}
\textbf{By the end of this chapter:} \emph{I can say to interfere, to insist, to put on record, to threaten and to hesitate.}

Complete the last lesson of chapter 451: I can say to interfere, to insist, to put on record, to threaten and to hesitate.
\end{chapteropening}

\section[\texorpdfstring{\ml{ഇടപെടുക} (iṭapeṭuka)}{ഇടപെടുക (iṭapeṭuka)}]{\ml{ഇടപെടുക} (iṭapeṭuka) — to interfere; to intervene}

\label{lesson:ML-C451-itapetuka}



\begin{quote}
Before the new one: say the Malayalam for to improve, then the Malayalam for to wrap up.
\end{quote}

\subsection*{You'll want to know: \ml{ഇടപെടുക}}
\textbf{\ml{ഇടപെടുക}} — \emph{iṭapeṭuka} — \textquotedblleft{}to interfere; to intervene\textquotedblright{}.

\begin{quote}
\ml{അവരുടെ} \ml{കാര്യത്തിൽ} \ml{ഞാൻ} \ml{ഇടപെട്ടില്ല}.
\end{quote}

\begin{quote}
\emph{avaruṭe kāryattil ñān iṭapeṭṭilla.}
\end{quote}

\begin{quote}
I did not interfere in their affairs.
\end{quote}

Past: \ml{ഇടപെട്ടു}, \emph{iṭapeṭṭu}.

In a matter: -\ml{ഇൽ} \ml{ഇടപെടുക}. The noun is \ml{ഇടപെടൽ}.

\begin{cousinweb}
\textbf{\ml{ഇടപെടുക}} (\emph{iṭapeṭuka}) is a native Malayalam compound: \ml{ഇട}, \textquotedblleft{}space, place\textquotedblright{} (DEDR 434), and \ml{പെടുക}, \textquotedblleft{}to be caught, to be in\textquotedblright{} (DEDR 3853): to step into the space between.
\end{cousinweb}

\begin{grammarlens}[title={What you've built}]
One more word for saying what happens if.
\end{grammarlens}

\subsection*{Your turn}
\begin{itemize}
\raggedright
  \item \emph{Say it:} \emph{iṭapeṭuka}
  \item \emph{Say it:} \emph{iṭapeṭuka}, once more
  \item \emph{Say it:} \emph{avaruṭe kāryattil ñān iṭapeṭṭilla.} — I did not interfere in their affairs.
  \item \emph{Recall:} say the Malayalam for to improve, then the Malayalam for to wrap up, then say \emph{iṭapeṭuka} again
\end{itemize}

\begin{tcolorbox}[breakable,colback=teal!4,colframe=teal!35!black,title={Before you move on}]
What does \ml{ഇടപെടുക} mean? (\textquotedblleft{}To interfere; to intervene\textquotedblright{}.) Say it once more.
\end{tcolorbox}
\section[\texorpdfstring{\ml{നിർബന്ധിക്കുക} (nirbandhikkuka)}{നിർബന്ധിക്കുക (nirbandhikkuka)}]{\ml{നിർബന്ധിക്കുക} (nirbandhikkuka) — to insist; to force}

\label{lesson:ML-C451-nirbandhikkuka}



\begin{quote}
Before the new one: say the Malayalam for to wrap up, then the Malayalam for to interfere.
\end{quote}

\subsection*{You'll want to know: \ml{നിർബന്ധിക്കുക}}
\textbf{\ml{നിർബന്ധിക്കുക}} — \emph{nirbandhikkuka} — \textquotedblleft{}to insist; to force\textquotedblright{}.

\begin{quote}
\ml{ഭക്ഷണം} \ml{കഴിക്കാൻ} \ml{അമ്മ} \ml{എന്നെ} \ml{നിർബന്ധിച്ചു}.
\end{quote}

\begin{quote}
\emph{bhakṣaṇaṁ kaḻikkān amma enne nirbandhiccu.}
\end{quote}

\begin{quote}
Mother insisted that I eat.
\end{quote}

Past: \ml{നിർബന്ധിച്ചു}, \emph{nirbandhiccu}.

The noun \ml{നിർബന്ധം}, insistence, is already known.

\begin{cousinweb}
\textbf{\ml{നിർബന്ധിക്കുക}} (\emph{nirbandhikkuka}) is built from Sanskrit \emph{nirbandha}, \textquotedblleft{}perseverance, insistence\textquotedblright{} (CDIAL 7355), with the Malayalam verb ending -\ml{ിക്കുക}. It is the verb of \ml{നിർബന്ധം}.
\end{cousinweb}

\begin{grammarlens}[title={What you've built}]
One more word for saying what happens if.
\end{grammarlens}

\subsection*{Your turn}
\begin{itemize}
\raggedright
  \item \emph{Say it:} \emph{nirbandhikkuka}
  \item \emph{Say it:} \emph{nirbandhikkuka}, once more
  \item \emph{Say it:} \emph{bhakṣaṇaṁ kaḻikkān amma enne nirbandhiccu.} — Mother insisted that I eat.
  \item \emph{Recall:} say the Malayalam for to wrap up, then the Malayalam for to interfere, then say \emph{nirbandhikkuka} again
\end{itemize}

\begin{tcolorbox}[breakable,colback=teal!4,colframe=teal!35!black,title={Before you move on}]
What does \ml{നിർബന്ധിക്കുക} mean? (\textquotedblleft{}To insist; to force\textquotedblright{}.) Say it once more.
\end{tcolorbox}
\section[\texorpdfstring{\ml{രേഖപ്പെടുത്തുക} (rēkhappeṭuttuka)}{രേഖപ്പെടുത്തുക (rēkhappeṭuttuka)}]{\ml{രേഖപ്പെടുത്തുക} (rēkhappeṭuttuka) — to put on record; to note down}

\label{lesson:ML-C451-rekhappetuttuka}



\begin{quote}
Before the new one: say the Malayalam for to interfere, then the Malayalam for to insist.
\end{quote}

\subsection*{You'll want to know: \ml{രേഖപ്പെടുത്തുക}}
\textbf{\ml{രേഖപ്പെടുത്തുക}} — \emph{rēkhappeṭuttuka} — \textquotedblleft{}to put on record; to note down\textquotedblright{}.

\begin{quote}
\ml{തൊഴിലാളികൾ} \ml{യോഗത്തിൽ} \ml{പ്രതിഷേധം} \ml{രേഖപ്പെടുത്തി}.
\end{quote}

\begin{quote}
\emph{toḻilāḷikaḷ yōgattil pratiṣēdhaṁ rēkhappeṭutti.}
\end{quote}

\begin{quote}
The workers put their protest on record at the meeting.
\end{quote}

Past: \ml{രേഖപ്പെടുത്തി}, \emph{rēkhappeṭutti}.

Also of votes and temperatures: \ml{വോട്ട്} \ml{രേഖപ്പെടുത്തി}, cast a vote.

\begin{cousinweb}
\textbf{\ml{രേഖപ്പെടുത്തുക}} (\emph{rēkhappeṭuttuka}) is built on \ml{രേഖ}, \textquotedblleft{}a line; a record\textquotedblright{}, from Sanskrit \emph{rēkhā}, \textquotedblleft{}a scratch, a line\textquotedblright{} (CDIAL 10810), with -\ml{പ്പെടുത്തുക} from native \ml{പെടുക} (DEDR 3853): to put a thing in writing.
\end{cousinweb}

\begin{grammarlens}[title={What you've built}]
One more word for saying what happens if.
\end{grammarlens}

\subsection*{Your turn}
\begin{itemize}
\raggedright
  \item \emph{Say it:} \emph{rēkhappeṭuttuka}
  \item \emph{Say it:} \emph{rēkhappeṭuttuka}, once more
  \item \emph{Say it:} \emph{toḻilāḷikaḷ yōgattil pratiṣēdhaṁ rēkhappeṭutti.} — The workers put their protest on record at the meeting.
  \item \emph{Recall:} say the Malayalam for to interfere, then the Malayalam for to insist, then say \emph{rēkhappeṭuttuka} again
\end{itemize}

\begin{tcolorbox}[breakable,colback=teal!4,colframe=teal!35!black,title={Before you move on}]
What does \ml{രേഖപ്പെടുത്തുക} mean? (\textquotedblleft{}To put on record; to note down\textquotedblright{}.) Say it once more.
\end{tcolorbox}
\section[\texorpdfstring{\ml{ഭീഷണിപ്പെടുത്തുക} (bhīṣaṇippeṭuttuka)}{ഭീഷണിപ്പെടുത്തുക (bhīṣaṇippeṭuttuka)}]{\ml{ഭീഷണിപ്പെടുത്തുക} (bhīṣaṇippeṭuttuka) — to threaten}

\label{lesson:ML-C451-bhisanippetuttuka}



\begin{quote}
Before the new one: say the Malayalam for to insist, then the Malayalam for to put on record.
\end{quote}

\subsection*{You'll want to know: \ml{ഭീഷണിപ്പെടുത്തുക}}
\textbf{\ml{ഭീഷണിപ്പെടുത്തുക}} — \emph{bhīṣaṇippeṭuttuka} — \textquotedblleft{}to threaten\textquotedblright{}.

\begin{quote}
\ml{കള്ളൻ} \ml{കത്തി} \ml{കാട്ടി} \ml{കടക്കാരനെ} \ml{ഭീഷണിപ്പെടുത്തി}.
\end{quote}

\begin{quote}
\emph{kaḷḷan katti kāṭṭi kaṭakkārane bhīṣaṇippeṭutti.}
\end{quote}

\begin{quote}
The thief threatened the shopkeeper with a knife.
\end{quote}

Past: \ml{ഭീഷണിപ്പെടുത്തി}, \emph{bhīṣaṇippeṭutti}.

The noun \ml{ഭീഷണി} is a threat.

\begin{cousinweb}
\textbf{\ml{ഭീഷണിപ്പെടുത്തുക}} (\emph{bhīṣaṇippeṭuttuka}) is built from Sanskrit \emph{bhīṣaṇa}, \textquotedblleft{}terrifying, frightening\textquotedblright{} (Monier-Williams), with -\ml{പ്പെടുത്തുക}, \textquotedblleft{}to cause to be\textquotedblright{}, from native \ml{പെടുക} (DEDR 3853).
\end{cousinweb}

\begin{grammarlens}[title={What you've built}]
One more word for saying what happens if.
\end{grammarlens}

\subsection*{Your turn}
\begin{itemize}
\raggedright
  \item \emph{Say it:} \emph{bhīṣaṇippeṭuttuka}
  \item \emph{Say it:} \emph{bhīṣaṇippeṭuttuka}, once more
  \item \emph{Say it:} \emph{kaḷḷan katti kāṭṭi kaṭakkārane bhīṣaṇippeṭutti.} — The thief threatened the shopkeeper with a knife.
  \item \emph{Recall:} say the Malayalam for to insist, then the Malayalam for to put on record, then say \emph{bhīṣaṇippeṭuttuka} again
\end{itemize}

\begin{tcolorbox}[breakable,colback=teal!4,colframe=teal!35!black,title={Before you move on}]
What does \ml{ഭീഷണിപ്പെടുത്തുക} mean? (\textquotedblleft{}To threaten\textquotedblright{}.) Say it once more.
\end{tcolorbox}
\section[\texorpdfstring{\ml{മടിക്കുക} (maṭikkuka)}{മടിക്കുക (maṭikkuka)}]{\ml{മടിക്കുക} (maṭikkuka) — to hesitate}

\label{lesson:ML-C451-matikkuka}



\begin{quote}
Before the new one: say the Malayalam for to put on record, then the Malayalam for to threaten.
\end{quote}

\subsection*{You'll want to know: \ml{മടിക്കുക}}
\textbf{\ml{മടിക്കുക}} — \emph{maṭikkuka} — \textquotedblleft{}to hesitate\textquotedblright{}.

\begin{quote}
\ml{സഹായം} \ml{ചോദിക്കാൻ} \ml{അവൾ} \ml{മടിച്ചു}.
\end{quote}

\begin{quote}
\emph{sahāyaṁ cōdikkān avaḷ maṭiccu.}
\end{quote}

\begin{quote}
She hesitated to ask for help.
\end{quote}

Past: \ml{മടിച്ചു}, \emph{maṭiccu}.

What one hesitates to do takes -\ml{ആൻ}. \ml{മടി} alone is laziness.

\begin{cousinweb}
\textbf{\ml{മടിക്കുക}} (\emph{maṭikkuka}) is native Dravidian: DEDR 4652 lists Malayalam \emph{maṭikka}, \textquotedblleft{}to grow tired, to be backward, averse, to doubt\textquotedblright{}.
\end{cousinweb}

\begin{grammarlens}[title={What you've built}]
One more word for saying what happens if.
\end{grammarlens}

\subsection*{Your turn}
\begin{itemize}
\raggedright
  \item \emph{Say it:} \emph{maṭikkuka}
  \item \emph{Say it:} \emph{maṭikkuka}, once more
  \item \emph{Say it:} \emph{sahāyaṁ cōdikkān avaḷ maṭiccu.} — She hesitated to ask for help.
  \item \emph{Recall:} say the Malayalam for to put on record, then the Malayalam for to threaten, then say \emph{maṭikkuka} again
\end{itemize}

\begin{tcolorbox}[breakable,colback=teal!4,colframe=teal!35!black,title={Before you move on}]
What does \ml{മടിക്കുക} mean? (\textquotedblleft{}To hesitate\textquotedblright{}.) Say it once more.
\end{tcolorbox}
